\section{Comparison and the scalar-curvature bound}\label{sec:barriers}

We combine the two reaction bounds with small Ricci curvature on
parabolic annuli. The inner comparison gives a spatial Ricci exponent
strictly below one; integrating its square through a scalar supersolution
then gives the required bounded scalar curvature.

For the inner Ricci comparison we will use a positive function of the form
\[
H=\delta^L(r^2+b_0\theta^2)^{-a/2},
\qquad a,L,b_0>0.
\]
Its radial diffusion has leading coefficient \(a(q-1)\), so we require
\(a>2\) to dominate the leading reaction coefficient \(2(q-1)\).
At the cap scale, the identity \(\delta^L=\theta^{L/\sigma}\) makes
its size proportional to \(\theta^{-(a-L/\sigma)}\). We will arrange
that this exponent lies between zero and one, allowing a bounded scalar
barrier to dominate the resulting Ricci-square source. The choices below also make the
parabolic-annulus estimate supply the lateral data for \(H\).

\subsection{Order of the choices}

We now choose one flow to prove Theorem~\ref{thm:main}. Set \(a=5/2\).
Since \(d\to2\) as \(q\to\infty\), choose \(q\) large enough that
\(2<d<a\), all universal thresholds above hold, and
\begin{equation}\label{eq:dimension-margin}
\mathfrak m:=(a-2)(q-1)-C\sqrt q-2a^2-1>0,
\end{equation}
where \(C\) is the universal constant in Proposition~\ref{prop:reaction}.
Next choose a sufficiently large even \(k\) so that, with
\begin{equation}\label{eq:barrier-exponents}
L=k-1,\qquad e=a-L/\sigma=a-d+d/k,
\end{equation}
we have \(0<e<1\), \(\nu>1\), and \(\sigma>1/2\).
These requirements are compatible because \(d>2\).
Define
\[
\xi=\min\{1/8,1/(2\sqrt L)\}.
\]
Choose the tolerances and cutoffs as in Section~\ref{sec:geometry}, and
fix the resulting flow. In particular \(C_*\) and \(D_1\) in
Proposition~\ref{prop:reaction} are now fixed finite numbers.
Only after this step choose \(b_0>0\) so that
\begin{equation}\label{eq:b0-choice}
b_0C_*<\mathfrak m.
\end{equation}
The remaining constants below are chosen for this same flow. Whenever a
later starting time is needed, we restrict to a later slice of the flow
already chosen; we do not repeat the construction or change its parameters.

In Figure~\ref{fig:scales}, \(\Gamma>\xi\) denotes a fixed outer
comparison radius whose size will be chosen in
\eqref{eq:exterior-choices}. The annular estimate is available for
every such fixed radius after moving the starting time later.

\begin{figure}[htbp]
\centering
\begin{tikzpicture}[x=1cm,y=1cm,>=Latex,font=\small]
  \fill[blue!9] (0,-.34) rectangle (2,.34);
  \fill[orange!12] (2,-.34) rectangle (6,.34);
  \fill[green!10] (6,-.34) rectangle (11,.34);
  \draw[->,thick] (0,0)--(11.35,0) node[right] {$r$};
  \foreach \x in {0,2,6,9} \draw (\x,-.11)--(\x,.11);
  \node[below=5pt] at (0,-.34) {$0$};
  \node[below=5pt] at (2,-.34) {$\theta$};
  \node[below=5pt] at (6,-.34) {$\xi\sqrt\delta$};
  \node[below=5pt] at (9,-.34) {$\Gamma\sqrt\delta$};
  \node[above] at (1,.42) {cap};
  \node[above,align=center] at (4,.42) {inner comparison};
  \node[above,align=center] at (8.5,.42) {annulus and exterior};
  \draw[decorate,decoration={brace,amplitude=4pt,mirror}] (0,-1.18)--(6,-1.18);
  \node[below,align=center,text width=6.2cm] at (3,-1.38)
     {$|\Ric|\le C\delta^L(r^2+b_0\theta^2)^{-a/2}$};
  \draw[decorate,decoration={brace,amplitude=4pt,mirror}] (6,-1.18)--(11,-1.18);
  \node[below,align=center,text width=4.7cm] at (8.5,-1.38)
     {$|\Ric|\le C$};
\end{tikzpicture}

\caption{The comparison regions at a late time, shown on a schematic
radius axis. The cap scale \(\theta\) is much smaller than the
parabolic scale \(\sqrt\delta\). The inner Ricci estimate bridges these
scales. Fixed parabolic-annulus control supplies lateral data for both
the inner and exterior comparisons; \(\Gamma\) is chosen later in
\eqref{eq:exterior-choices}.}
\label{fig:scales}
\end{figure}

\subsection{The inner Ricci estimate}

\begin{proposition}\label{prop:ricci-barrier}
There is a constant \(C\) such that, at all sufficiently late times,
\begin{equation}\label{eq:inner-ricci}
|\Ric|\le C\delta^L(r^2+b_0\theta^2)^{-a/2}
\qquad(0\le r\le\xi\sqrt\delta).
\end{equation}
In particular,
\begin{equation}\label{eq:ricci-power}
|\Ric|\le Cr^{-e}\qquad(0<r\le\xi\sqrt\delta).
\end{equation}
\end{proposition}

\begin{proof}
Put \(D_0=r^2+b_0\theta^2\) and \(H=\delta^LD_0^{-a/2}\).
This is a positive smooth geometric function on both closed caps,
including their pole orbits, because \(r^2\) is smooth there.
Using \(\theta_t=-\sigma\theta/\delta\) and
Equation~\eqref{eq:operator}, direct differentiation gives
\begin{equation}\label{eq:inner-barrier-computation}
\begin{split}
\frac{\D H}{H}
&=-\frac L\delta+\frac{a\sigma b_0\theta^2}{\delta D_0}
 +\frac{a(q-1)+2az}{D_0}-\frac{a(a+2)zr^2}{D_0^2}\\
&\ge-\frac L\delta+\frac{a(q-1)-a^2z}{D_0}.
\end{split}
\end{equation}
The first line has a nonnegative extra time term, and \(r^2\le D_0\)
gives the second. The identities extend to a pole by continuity; in
particular the apparent \(r^{-1}H_r\) term has a finite limit.

In the indicated region \(z\le2\) for large \(q\), and
\[
LD_0/\delta\le L\xi^2+Lb_0\delta^{2\sigma-1}\le1
\]
at sufficiently late times. Where \(|\Ric|>0\) and \(r>0\), the two
bounds in~\eqref{eq:reaction-two} imply
\[
cD_0=cr^2+b_0c\theta^2
\le2(q-1)+C\sqrt q+b_0C_*.
\]
Both inequalities may be added even if \(c\) is negative. At a pole,
the second reaction bound alone implies the same upper bound for
\(cD_0\). Therefore
\begin{equation}\label{eq:strict-ricci-super}
\frac{\D H}{H}-c
\ge\frac{-1+a(q-1)-2a^2-2(q-1)-C\sqrt q-b_0C_*}{D_0}>0
\end{equation}
by~\eqref{eq:dimension-margin} and~\eqref{eq:b0-choice}.

On the moving lateral boundary \(r=\xi\sqrt\delta\),
\[
H=\delta^{L-a/2}
    (\xi^2+b_0\delta^{2\sigma-1})^{-a/2}
\asymp\delta^{L-a/2}.
\]
Lemma~\ref{lem:annuli} and \(\nu=k-d/2\) give
\begin{equation}\label{eq:lateral-ratio}
\frac{|\Ric|}{H}\le C\delta^{\nu-1-(L-a/2)}
=C\delta^{(a-d)/2}.
\end{equation}
This ratio is uniformly bounded at late times. Choose a sufficiently
late fixed starting time \(t_0\). On that compact starting slice, \(H\)
has a positive minimum on the caps, so a single multiple of \(H\)
strictly dominates both initial and lateral data for every subsequent
compact time slab.

At a first contact of \(|\Ric|\) with this positive barrier, the Ricci
norm is nonzero and hence smooth nearby. Equation~\eqref{eq:ricci-norm}
and the strict inequality~\eqref{eq:strict-ricci-super} rule out the
contact by the maximum principle. This argument also applies at a pole,
which is interior to the smooth manifold. It proves~\eqref{eq:inner-ricci}
with one multiplier for all times \(t_0\le t<T\).

Finally \(\delta^L=\theta^{L/\sigma}=\theta^{a-e}\), with \(a-e>0\).
For \(r\ge\theta\),
\[
H\le\theta^{a-e}r^{-a}\le r^{-e}.
\]
For \(0<r\le\theta\),
\[
H\le b_0^{-a/2}\theta^{-e}\le b_0^{-a/2}r^{-e}.
\]
These prove~\eqref{eq:ricci-power}.
\end{proof}

\subsection{A scalar barrier at the pole}

\begin{proposition}\label{prop:scalar-inner}
The scalar curvature is uniformly bounded above on
\(0\le r\le\xi\sqrt\delta\) at late times.
\end{proposition}

\begin{proof}
For a real exponent \(\alpha\), Equation~\eqref{eq:operator} gives
\begin{equation}\label{eq:power-operator}
\D(r^\alpha)=-\alpha(q-1+\alpha z)r^{\alpha-2}
\qquad(r>0).
\end{equation}
Put \(\ell_0=2-2e\in(0,2)\). In particular,
\begin{equation}\label{eq:scalar-powers}
\D(-r^{\ell_0})=\ell_0(q-1+\ell_0z)r^{-2e},
\qquad \D(r^{-1})=(q-1-z)r^{-3}.
\end{equation}
Write \(|\Ric|\le K_Rr^{-e}\) as in Proposition~\ref{prop:ricci-barrier}.
Choose
\[
C_2\ge\frac{2K_R^2}{\ell_0(q-1)}.
\]
Since \(0\le z\le2<q-1\), every function
\begin{equation}\label{eq:scalar-barrier}
S_\varepsilon=C_1-C_2r^{\ell_0}+\varepsilon r^{-1},
\qquad \varepsilon>0,
\end{equation}
satisfies \(\D S_\varepsilon\ge2|\Ric|^2\) at positive radii in the
inner region.

Choose one late starting time \(t_0\) and let
\(\rho_0=\xi\sqrt{T-t_0}\). Smoothness bounds scalar curvature on
the starting slice by a constant \(K_0\). On the moving lateral boundary,
Lemma~\ref{lem:annuli}, \(\nu>1\), and
\(|R|\le\sqrt{m+1}|\Ric|\) give a uniform upper bound \(K_\partial\).
The choice
\begin{equation}\label{eq:C1-choice}
C_1>\max\{K_0,K_\partial\}+C_2\rho_0^{\ell_0}
\end{equation}
therefore dominates the initial and lateral scalar data for every
\(\varepsilon>0\).

Fix \(\varepsilon>0\) and \(t_1<T\). Scalar curvature is bounded on
the compact slab \([t_0,t_1]\). For sufficiently small \(\eta>0\), the
level \(r=\eta\) lies in the regular inner region throughout that slab
and \(\varepsilon/\eta\) makes \(S_\varepsilon\) dominate its scalar
data. Apply comparison to \(R-S_\varepsilon\) on
\[
\{(x,t):t_0\le t\le t_1,\ \eta\le r(x,t)\le\xi\sqrt{T-t}\}.
\]
The moving outer boundary is regular in the monotone inner region.
Both initial and lateral boundaries have been controlled, so the usual
interior first-maximum argument applies. Let \(\eta\downarrow0\).
This proves \(R\le S_\varepsilon\) at every positive radius on the slab.

Only the auxiliary excision radius depends on \(\varepsilon,t_1\);
the constants in~\eqref{eq:C1-choice} do not. Letting the slab range up to
\(T\) and then sending \(\varepsilon\downarrow0\) pointwise gives
\[
R\le C_1-C_2r^{\ell_0}\le C_1\qquad(r>0).
\]
At every fixed time \(t<T\), continuity gives the same bound at the
pole orbits. No smoothness of \(r^{\ell_0}\) at the pole or uniform
regularity at \(t=T\) was used.
\end{proof}

\subsection{The exterior and the maximal time}

\begin{proposition}\label{prop:exterior}
There is a constant \(C\) such that
\(|\Ric|\le C\) on \(r\ge\xi\sqrt\delta\) at late times.
\end{proposition}

\begin{proof}
Choose \(b_1>2D_1+1\), where \(D_1\) is the global reaction constant
in Proposition~\ref{prop:reaction}, and choose
\begin{equation}\label{eq:exterior-choices}
\Gamma>\xi,\qquad \Gamma^2\ge\max\{2b_1,4(q-1)\}.
\end{equation}
On the entire geometric exterior \(r\ge\Gamma\sqrt\delta\), put
\(H_o=1-b_1\delta/r^2\). This is smooth across the equator and satisfies
\(1/2\le H_o\le1\). Equation~\eqref{eq:operator} gives
\begin{equation}\label{eq:outer-barrier}
\D H_o=\frac{b_1}{r^2}
\left[1-\bigl(2(q-1)-4z\bigr)\frac{\delta}{r^2}\right]
\ge\frac{b_1}{2r^2}>\frac{D_1H_o}{r^2}.
\end{equation}
Here only \(z\ge0\) was needed; no global smallness of \(z-1\) is
required.

The fixed parabolic radius \(\Gamma\) eventually lies in the region of
Lemma~\ref{lem:annuli}. The annular estimate and \(\nu>1\) bound \(|\Ric|\) on
the moving lateral boundary uniformly. On a sufficiently late starting
slice, compactness and \(H_o\ge1/2\) allow one multiplier to dominate
the initial data. The same first-contact argument as in
Proposition~\ref{prop:ricci-barrier} proves \(|\Ric|\le CH_o\le C\)
throughout the exterior. The equator is interior, since both hemispheres
are included. On the remaining fixed annulus
\(\xi\le r/\sqrt\delta\le\Gamma\), use Lemma~\ref{lem:annuli}
directly.
\end{proof}

\begin{proof}[Proof of Theorem~\ref{thm:main}]
Take the single flow selected at the start of Section~\ref{sec:barriers}.
Propositions~\ref{prop:scalar-inner} and~\ref{prop:exterior} give a
uniform scalar upper bound on all of \(M\) for sufficiently late times.
There are only finitely many required starting-time restrictions, and
every parameter and constant has already been fixed, so a common such
time exists. Earlier times form a compact smooth interval and contribute
a finite bound. Lemma~\ref{lem:gradients} supplies the scalar lower bound.
Thus \(\sup_{M\times[0,T)}|R|<\infty\).

By Lemma~\ref{lem:pole-size}, \(\phi(0,t)\le C\theta(t)\to0\).
At a pole orbit, smoothness gives \(\phi_s=0\). Since \(p=2\), a
two-plane tangent to the \(\Sp^p\) orbit has sectional curvature
\[
j(0,t)=\phi(0,t)^{-2}\longrightarrow\infty.
\]
Consequently \(\max_M|\Rm|\to\infty\). A smooth extension on the same
compact manifold would bound curvature on a compact time neighborhood
of \(T\), which is impossible. The constructed flow exists smoothly for
every \(t<T\); closed-manifold uniqueness therefore identifies \(T\)
with the maximal existence time for its smooth initial metric. The product
\(\Sp^2\times\Sp^{q+1}\) is closed, connected, and has dimension
\(q+3\ge4\), as required.
\end{proof}

\begin{corollary}[Curvature rates in a fixed dimension]\label{cor:rates}
For every sufficiently large integer \(q\), set
\[
d_q=\frac{q+1-\sqrt{(q+1)(q-7)}}2.
\]
For every sufficiently large even integer \(k\), with the threshold
allowed to depend on \(q\), there is a smooth Ricci flow on
\(\Sp^2\times\Sp^{q+1}\) with finite maximal time \(T\), uniformly
bounded scalar curvature, and constants \(0<c\le C<\infty\) such that
\[
c(T-t)^{-2k/d_q}\le\max_M|\Rm|(\cdot,t)
\le C(T-t)^{-2k/d_q}
\]
for all sufficiently late times. In particular, one fixed dimension
admits such flows with an unbounded discrete set of power-law Type-II
curvature exponents. The flow, \(T\), and the constants may depend on
\(q\) and \(k\).
\end{corollary}

\begin{proof}
Every threshold used to choose \(q\) at the beginning of this section
is independent of \(k\). Once such a \(q\) is fixed, all sufficiently
large even \(k\) satisfy the requirements in
\eqref{eq:barrier-exponents}; Proposition~\ref{prop:input} then permits
the remaining choices. The preceding comparisons apply to each resulting
flow, with constants that may depend on it.

Here \(d=d_q\) and \(\theta=(T-t)^{k/d_q}\).
Proposition~\ref{prop:full-curvature} gives
\(\max_M|\Rm|\le C\theta^{-2}\).
At either pole orbit, Lemma~\ref{lem:pole-size} gives
\(\phi(0,t)\le C\theta\), and the tangential sectional curvature is
\(j(0,t)=\phi(0,t)^{-2}\). Hence
\(\max_M|\Rm|\ge c\theta^{-2}\).
These are the claimed two-sided bounds. Since \(2k/d_q>1\) and tends
to infinity with \(k\), the flows are Type~II and their exponents are
unbounded with \(q\) fixed.
\end{proof}
